\documentclass[11pt,a4paper]{article}
\usepackage[utf8]{inputenc}
\usepackage[T1]{fontenc}
\usepackage[german,english]{babel}
\usepackage{amsmath,amssymb,amsthm}
\usepackage{geometry}
\usepackage{booktabs}
\usepackage{hyperref}

\geometry{verbose,tmargin=2.5cm,bmargin=2.5cm,lmargin=2.5cm,rmargin=2.5cm}

\newtheorem{definition}{Definition}
\newtheorem{proposition}{Proposition}
\newtheorem{theorem}{Satz}

\title{\textbf{Radiale Quotiententopologie und meromorphe Struktur der Epstein-Zetafunktion des $E_8$-Gitters}}
\author{\textbf{Thomas Hoffbauer} \\ \small Computer Institut Bamberg (CIB)}
\date{1. August 2026}

\begin{document}
	
	\maketitle
	
	\begin{abstract}
		\selectlanguage{english}
		We construct the radial quotient topology $\tau_{\text{rad}}$ on the $E_8$ root lattice generated by the norm quadratic form $Q(x) = \|x\|^2/2 \in \mathbb{N}_0$. We prove that the Kolmogorov quotient space $E_8 / \sim$ is homeomorphic to the discrete set of non-negative integers $\mathbb{N}_0$. Furthermore, we evaluate the associated radial Epstein zeta function $Z_{E_8}(s)$ via the theta series $\Theta_{E_8}(q)$. We establish that $Z_{E_8}(s)$ possesses an analytic continuation with exactly one simple pole at $s=8$ with residue $\text{Res}_{s=8} Z_{E_8}(s) = \pi^4 / 3$, while the apparent singularity at $s=2$ is removable with $Z_{E_8}(2) = -30\zeta(3)/\pi^2$. Finally, we derive the exact asymptotic ball growth $|K_N| \sim \frac{2\pi^4}{3} N^4$ and the arithmetic Möbius inversion for primitive lattice vectors.
	\end{abstract}
	
	\selectlanguage{german}
	
	\section{Einleitung und Grundlegende Definitionen}
	
	Das $E_8$-Wurzelgitter nimmt aufgrund seiner herausragenden Symmetrie- und Selbstdualitätseigenschaften eine zentrale Stellung in der diskreten Geometrie und der Theorie der Modulformen ein. In dieser Arbeit untersuchen wir die arithmetische Schichtenarchitektur von $\mathbf{E}_8$ unter einer natürlichen radialen Quotiententopologie.
	
	Sei $\mathbf{E}_8 \subset \mathbb{R}^8$ das gerade, selbstduale $E_8$-Gitter. Wir betrachten die normierte quadratische Form $Q(x) := \frac{\|x\|^2}{2} \in \mathbb{N}_0$ für alle $x \in \mathbf{E}_8$.
	
	\begin{definition}[Schalen-Partition]
		Für $n \in \mathbb{N}_0$ definieren wir die $n$-te $E_8$-Normschale $S_n$ als das Urbild der Menge $\{n\}$ unter der quadratischen Form $Q$:
		$$S_n := Q^{-1}(\{n\}) = \left\{ x \in \mathbf{E}_8 \;\middle|\; \frac{\|x\|^2}{2} = n \right\}$$
		Der Gesamtraum zerfällt in die disjunkte Vereinigung $\mathbf{E}_8 = \bigsqcup_{n=0}^{\infty} S_n$.
	\end{definition}
	
	\begin{definition}[Radiale Quotiententopologie]
		Die radiale Quotiententopologie $\tau_{\text{rad}}$ auf $\mathbf{E}_8$ ist die Initialtopologie bezüglich der Abbildung $Q: \mathbf{E}_8 \to \mathbb{N}_0$, wobei $\mathbb{N}_0$ die diskrete Topologie trägt:
		$$\tau_{\text{rad}} := \left\{ Q^{-1}(V) \;\middle|\; V \subseteq \mathbb{N}_0 \right\} = \left\{ \bigcup_{n \in V} S_n \;\middle|\; V \subseteq \mathbb{N}_0 \right\}$$
	\end{definition}
	
	---
	
	\section{Topologische und Ordnungstheoretische Eigenschaften}
	
	\begin{proposition}[Quotienten-Struktur]
		Der topologische Raum $(\mathbf{E}_8, \tau_{\text{rad}})$ ist ein pseudometrischer Raum unter der Pseudometrik $d_{\text{rad}}(x,y) = |Q(x) - Q(y)|$. Der Kolmogorov-Quotient $X / \sim$ bezüglich der Äquivalenzrelation $x \sim y \iff Q(x) = Q(y)$ ist homöomorph zur diskreten Menge $\mathbb{N}_0$. Jede Schale $S_n$ ist invariant unter der Automorphismengruppe $\operatorname{Aut}(\mathbf{E}_8)$ und zerfällt in endlich viele Orbits.
	\end{proposition}
	
	\begin{proposition}[Alexandrov-Radialtopologie und Kompaktheit]
		Die durch die Präordnung $x \preceq y \iff Q(x) \le Q(y)$ induzierte Alexandrov-Topologie $\tau_{\uparrow}$ besteht aus den Urbildern aufwärts abgeschlossener Teilmengen von $\mathbb{N}_0$. In $(\mathbf{E}_8, \tau_{\uparrow})$ sind die inneren Vollkugeln $K_N := \{x \in \mathbf{E}_8 \mid Q(x) \le N\}$ abgeschlossen und als endliche Teilmengen kompakt.
	\end{proposition}
	
	---
	
	\section{Arithmetik der Schalen und Asymptotik}
	
	Die Kardinalitäten $|S_n|$ der Schalen werden direkt durch die Theta-Reihe des $E_8$-Gitters bestimmt.
	
	\begin{theorem}[Jacobi-Theta-Identität und Ballasymptotik]
		Für die Erzeugendenfunktion der Schalenpunktzahlen $|S_n|$ gilt:
		$$\Theta_{\mathbf{E}_8}(q) = \sum_{x \in \mathbf{E}_8} q^{Q(x)} = 1 + 240 \sum_{n=1}^{\infty} \sigma_3(n) q^n \implies |S_n| = 240 \, \sigma_3(n) \quad (n \ge 1)$$
		wobei $\sigma_3(n) = \sum_{d|n} d^3$ die Teilerfunktion dritter Potenz bezeichnet. Für das kumulierte Punktvolumen der inneren Vollkugeln $K_N$ gilt das asymptotische Wachstumsgesetz:
		$$|K_N| = 1 + 240 \sum_{n=1}^{N} \sigma_3(n) = \frac{2\pi^4}{3} N^4 + \mathcal{O}(N^3 \log N)$$
	\end{theorem}
	
	---
	
	\section{Analytische Fortsetzung der radialen Dirichlet-Reihe}
	
	\begin{theorem}[Meromorphe Struktur von $Z_{\mathbf{E}_8}(s)$]
		Für $\operatorname{Re}(s) > 8$ ist die radiale Epstein-Zetafunktion des $E_8$-Gitters definiert durch:
		$$Z_{\mathbf{E}_8}(s) := \sum_{x \in \mathbf{E}_8 \setminus \{0\}} \|x\|^{-s} = 240 \cdot 2^{-s/2} \cdot \zeta\left(\frac{s}{2}\right) \cdot \zeta\left(\frac{s}{2} - 3\right)$$
		$Z_{\mathbf{E}_8}(s)$ besitzt eine meromorphe Fortsetzung auf $\mathbb{C}$ mit exakt einer Polstelle erster Ordnung bei $s = 8$. Das Residuum beträgt:
		$$\operatorname*{Res}_{s=8} Z_{\mathbf{E}_8}(s) = \frac{\pi^4}{3}$$
		Die Stelle $s = 2$ ist eine hebbare Singularität mit dem regulären Funktionswert:
		$$Z_{\mathbf{E}_8}(2) = -\frac{30\,\zeta(3)}{\pi^2}$$
	\end{theorem}
	
	\begin{proof}
		Die Produktdarstellung folgt direkt durch Einsetzen von $|S_n| = 240 \sigma_3(n)$ in die Dirichlet-Reihe und die klassische Identität $\sum_{n=1}^\infty \sigma_3(n) n^{-u} = \zeta(u)\zeta(u-3)$ mit $u = s/2$.
		Für $s \to 8$ ($u \to 4$) besitzt $\zeta(u-3) = \zeta(u-3)$ bei $u=1$ den Hauptteil $\frac{1}{u-1} = \frac{2}{s-8}$. Das Residuum berechnet sich somit zu $240 \cdot 2^{-4} \cdot \zeta(4) \cdot 2 = 30 \cdot \frac{\pi^4}{90} = \frac{\pi^4}{3}$.
		An der Stelle $s = 2$ ($u = 1$) besitzt $\zeta(u)$ einen einfachen Pol mit Hauptteil $\frac{1}{u-1}$, während $\zeta(u-3) = \zeta(-2)$ eine einfache triviale Nullstelle aufweist. Pol und Nullstelle heben sich auf. Über den Limes erhält man $Z_{\mathbf{E}_8}(2) = 240 \cdot 2^{-1} \cdot \zeta'(-2)$. Mit der bekannten Identität $\zeta'(-2) = -\frac{\zeta(3)}{4\pi^2}$ folgt der Wert $Z_{\mathbf{E}_8}(2) = -\frac{30\zeta(3)}{\pi^2}$.
	\end{proof}
	
	\begin{theorem}[Arithmetische Möbius-Inversion für primitive Vektoren]
		Sei $\mathbf{E}_8^{\text{prim}} \subset \mathbf{E}_8$ die Menge der primitiven Gittervektoren (nicht durch ein Ganzes $m \ge 2$ teilbar). Dann gilt im Konvergenzbereich:
		$$Z_{\mathbf{E}_8}^{\text{prim}}(s) := \sum_{x \in \mathbf{E}_8^{\text{prim}}} \|x\|^{-s} = \frac{Z_{\mathbf{E}_8}(s)}{\zeta(s)} = 240 \cdot 2^{-s/2} \cdot \frac{\zeta(s/2)\,\zeta(s/2 - 3)}{\zeta(s)}$$
	\end{theorem}
	
	\begin{proof}
		Jeder Vektor $x \in \mathbf{E}_8 \setminus \{0\}$ lässt sich eindeutig als $x = m \cdot y$ schreiben mit $m \in \mathbb{N}$ und $y \in \mathbf{E}_8^{\text{prim}}$. Wegen $\|x\|^{-s} = m^{-s} \|y\|^{-s}$ folgt $Z_{\mathbf{E}_8}(s) = \zeta(s) Z_{\mathbf{E}_8}^{\text{prim}}(s)$.
	\end{proof}
	
	\section{Fazit}
	Die radiale Quotiententopologie $(\mathbf{E}_8, \tau_{\text{rad}})$ liefert einen streng rigorosen, strukturellen Rahmen zur Analyse der Schalenarithmetik des $E_8$-Gitters. Die analytischen Eigenschaften der zugehörigen Epstein-Zetafunktion spiegeln die Volumeneigenschaften und die Verteilung primitiver Vektoren ohne verbleibende Singularitäten abseits des Hauptpols bei $s=8$ wider.
	
\end{document}